\section*{Reading notes}
\subsection*{\texttt{GELFAND-P001435}}
\par\noindent\leavevmode\phantomsection\label{GELFAND-0156-U001}\label{GELFAND-P001435:heading-date}\textbf{\texttt{GELFAND-0156-U001}} (heading-date).\nopagebreak[4]\par
At the upper left is 22.2.8…; the final digit has overlapping or intersecting strokes.\par
Candidates: 22.2.88; 22.2.82. Unresolved in the diplomatic source.\par
Historical source review: Original PNG, bbox [170,45,345,125], original-pixel coordinates (x0,y0,x1,y1) from the upper-left corner; crop [125,30,385,145] inspected at 3×. The catalogue date 1988-02-22 was not used to settle the final digit silently.\par
Current source comparison: The final digit remains uncertain after original-pixel 3×, 5× and 10× inspection; the catalogue year cannot decide the strokes.\par\medskip
\subsection*{\texttt{GELFAND-P001436}}
\subsection*{\texttt{GELFAND-P001437}}
\subsection*{\texttt{GELFAND-P001438}}
\par\noindent\leavevmode\phantomsection\label{GELFAND-0156-U002}\label{GELFAND-P001438:operator-coefficient}\textbf{\texttt{GELFAND-0156-U002}} (operator-coefficient).\nopagebreak[4]\par
In (1), the subscript after u appears to be N−1 or N−i. The last stroke has a sloping upper stroke and a lower foot; no separate dot of an i can be distinguished.\par
Candidates: u\_\{N-1\}(z); u\_\{N-i\}(z). Unresolved in the diplomatic source.\par
Historical source review: Original coordinates [827,552,944,645]; repeat crop [837,550,966,660] at 10×. The more literal N−1 is retained in the formula. Mathematical check: with the subscript N−i the coefficients would vary with the summation index; this does not justify correcting the visible N−1.\par
Current source comparison: No independent dot establishes i in the new 3×/5×/10× views. The literal N−1 remains; the separate conditional N=2 comparison is an editorial consistency check only.\par\medskip
\par\noindent\leavevmode\phantomsection\label{GELFAND-0156-U003}\label{GELFAND-P001438:local-functional-argument}\textbf{\texttt{GELFAND-0156-U003}} (local-functional-argument).\nopagebreak[4]\par
In integral (*), the middle argument is written u(…); the inner sign joins closely to the parenthesis and resembles x or z.\par
Candidates: f(x,u(x),u'(x)); f(x,u(z),u'(x)). Unresolved in the diplomatic source.\par
Historical source review: bbox [692,2202,842,2316], formula crop [459,2160,1012,2348] at 3×. Candidate x is shown in the main text. Agreement with the adjacent arguments and dx was not treated as proof of the reading.\par
Current source comparison: The whole formula and the tighter middle-u argument were inspected at original-pixel 3×/5×/10×. The close joins do not establish x rather than z; locality reasoning does not decipher the sign.\par\medskip
\par\noindent\leavevmode\phantomsection\label{GELFAND-0156-U004}\label{GELFAND-P001438:locality-diagram}\textbf{\texttt{GELFAND-0156-U004}} (locality-diagram).\nopagebreak[4]\par
To the left of the locality condition, a sloping line runs from the lower left to the upper right. Two small protrusions point upwards; upper or left strokes are broken in places. No labels or arrowheads are visible.\par
Candidates: none recorded. Unresolved in the diplomatic source.\par
Historical source review: bbox [368,2500,850,2670], inspected at 3×. It remains necessary to check whether the breaks belong to the drawing or represent strokes lost in scanning. The geometry was not replaced by an invented smooth curve; SourceFigure was retained.\par
Current source comparison: The historical SourceFigure decision is preserved above. The new supplement adds the exact original crop and a partial native trace of the clear broken strokes. Separate 3×/5×/10× views of both breaks do not determine their cause; no joins, axes, values or arrows are invented. PDF layout remains unreviewed.\par\medskip
\par\noindent\leavevmode\phantomsection\label{GELFAND-0156-U005}\label{GELFAND-P001438:locality-cancelled}\textbf{\texttt{GELFAND-0156-U005}} (locality-cancelled).\nopagebreak[4]\par
Between “sum” and “increments”, two almost vertical sloping strokes are heavily crossed out.\par
Candidates: 2; beginning of “two” [source: «двух»]. Unresolved in the diplomatic source.\par
Historical source review: bbox [1885,2625,1970,2720], crop [1790,2615,1960,2722] at 5×. The original letters cannot confidently be separated from the cancellation.\par
Current source comparison: New original-pixel 3×/5×/10× views do not separate the underlying strokes; both candidates remain cancelled and uncertain.\par\medskip
\par\noindent\leavevmode\phantomsection\label{GELFAND-0156-U006}\label{GELFAND-P001438:general-functional-phrase}\textbf{\texttt{GELFAND-0156-U006}} (general-functional-phrase).\nopagebreak[4]\par
Below “(locality!)”, a short phrase is underlined. Its beginning resembles “5-й” / “Б-й”; next is an abbreviation resembling “возм.”, a word with two tall loops read as “обобщ.”, and “ф-й:” / “ф-и:”.\par
Candidates: 5th poss. generaliz. of f-ns: [source: «5-й возм. обобщ. ф-й:»]; B-th poss. generaliz. of f-ns: [source: «Б-й возм. обобщ. ф-й:»]; 5th kind. general. of f-ns: [source: «5-й вид. общ. ф-й:»]; B-th poss. generaliz. of f-n: [source: «Б-й возм. обобщ. ф-и:»]. Unresolved in the diplomatic source.\par
Historical source review: bbox [305,2820,1085,2942]; crops [310,2820,1102,2954] at 3× and [295,2810,731,2950] at 4×. The opening phrase is unresolved; its grammatical and semantic difficulty was not removed by conjecturing “if one takes a general functional”. Rechecking the general-opening and common-functional crops favoured “возм.” and “обобщ.” over the earlier “вид. общ.”. The phrase was not replaced with a presumed expanded argument.\par
Current source comparison: The whole phrase at 3× and its opening, middle and ending separately at 5× and 10× retain the same doubts. The English abbreviations are candidate translations, not an expansion of the source into an established theorem.\par\medskip
\par\noindent\leavevmode\phantomsection\label{GELFAND-0156-U007}\label{GELFAND-P001438:general-functional-cancelled}\textbf{\texttt{GELFAND-0156-U007}} (general-functional-cancelled).\nopagebreak[4]\par
Between “then” and “functional [abbr.]”, a short word is crossed out; its original strokes resemble от or он.\par
Candidates: from [source: «от»]; it/he [source: «он»]. Unresolved in the diplomatic source.\par
Historical source review: bbox [785,3006,894,3072], crop [780,2990,916,3088] at 6×. The cancellation is retained descriptively.\par
Current source comparison: The 3×/5×/10× views do not settle the cancelled word; the sentence was not used to supply it.\par\medskip
\par\noindent\leavevmode\phantomsection\label{GELFAND-0156-U008}\label{GELFAND-P001438:type-zero-cancelled}\textbf{\texttt{GELFAND-0156-U008}} (type-zero-cancelled).\nopagebreak[4]\par
Before the final phrase “This is type 0” is an isolated, heavily crossed-out sign resembling 1 or the Cyrillic letter И.\par
Candidates: 1; I [Cyrillic letter «И»]. Unresolved in the diplomatic source.\par
Historical source review: bbox [393,3120,457,3195], crop [365,3110,480,3200] at 6×. It was not treated as an additional item number; it is marked as a crossed-out sign.\par
Current source comparison: The 3×/5×/10× views still show an isolated cancelled sign; no extra numbered item is introduced.\par\medskip
\subsection*{\texttt{GELFAND-P001439}}
\par\noindent\leavevmode\phantomsection\label{GELFAND-0156-U031}\label{GELFAND-P001439:analysis-verb}\textbf{\texttt{GELFAND-0156-U031}} (analysis-verb).\nopagebreak[4]\par
In the first sentence after “it”, the verb begins “форми…” / “формули…”; next is a tall loop resembling в and an ending resembling “…ался”.\par
Candidates: was taking shape [source: «формировался»]; was being formulated [source: «формулировался»]; is taking shape [source: «формируется»]. Unresolved in the diplomatic source.\par
Historical source review: Original 2408×3480; bbox (1542,489,2100,607), magnification 4×. On rechecking, “формируется” ceased to be a confident reading; past tense is not asserted solely from the syntax of the following line.\par
Current source comparison: The full verb at 3× and its ending at 5×/10× do not establish the tense or all middle letters; the English line keeps all three possibilities.\par\medskip
\par\noindent\leavevmode\phantomsection\label{GELFAND-0156-U009}\label{GELFAND-P001439:rationalize-cancelled}\textbf{\texttt{GELFAND-0156-U009}} (rationalize-cancelled).\nopagebreak[4]\par
After “when everyone/all”, a short word ending in a rounded sign is crossed out.\par
Candidates: this [source: «это»]. Unresolved in the diplomatic source.\par
Historical source review: bbox [567,606,726,669], line crop [225,582,1710,707] at 3×. “Это” is likely, but horizontal strokes cover the middle. The uncancelled “пытаешь” was left without adding the reflexive suffix -ся.\par
Current source comparison: The 3×/5×/10× detail leaves the cancelled middle letters obscured. No reflexive suffix is supplied to the neighbouring verb.\par\medskip
\par\noindent\leavevmode\phantomsection\label{GELFAND-0156-U032}\label{GELFAND-P001439:rationalize-verb}\textbf{\texttt{GELFAND-0156-U032}} (rationalize-verb).\nopagebreak[4]\par
In the line “when everyone/all … to rationalize”, the middle uncancelled word starts “пыта…” and ends with a crowded group of strokes.\par
Candidates: tried [source: «пытались»]; you try [source: «пытаешь», without a supplied reflexive suffix]. Unresolved in the diplomatic source.\par
Historical source review: Original 2408×3480; bbox (225,582,1710,707), magnification 3×. Rechecking permits “пытались” instead of the initial “пытаешь”; the ending remains open and was not corrected to match grammatical expectations.\par
Current source comparison: The verb at 3×/5×/10× still does not establish its ending; both historical candidates are retained without grammatical repair.\par\medskip
\par\noindent\leavevmode\phantomsection\label{GELFAND-0156-U010}\label{GELFAND-P001439:continuity-cancelled}\textbf{\texttt{GELFAND-0156-U010}} (continuity-cancelled).\nopagebreak[4]\par
After the isolated “And”, a group of letters with a descending final stroke is crossed out; it resembles “spectrum”.\par
Candidates: spectrum [source: «спектр»]. Unresolved in the diplomatic source.\par
Historical source review: bbox [404,687,644,767], crop [225,681,2190,846] at 3×. The original letters could not be completely separated from the crossing-out lines.\par
Current source comparison: New 3×/5×/10× inspection does not separate the cancellation from every letter. “Spectrum” remains a candidate only.\par\medskip
\par\noindent\leavevmode\phantomsection\label{GELFAND-0156-U012}\label{GELFAND-P001439:type-two-cancelled}\textbf{\texttt{GELFAND-0156-U012}} (type-two-cancelled).\nopagebreak[4]\par
Between “functional [abbr.]” and “of type 2”, two or three short signs are crossed out.\par
Candidates: deg. [source: «ст.»]; from [source: «от»]. Unresolved in the diplomatic source.\par
Historical source review: bbox [377,2006,485,2087], crop [390,2005,502,2091] at 6×. The possible abbreviation “ст.” agrees with the previous line, but is only a candidate, not a resolution.\par
Current source comparison: The 3×/5×/10× views do not settle the two or three signs. The preceding definition supplies no palaeographic proof.\par\medskip
\par\noindent\leavevmode\phantomsection\label{GELFAND-0156-U013}\label{GELFAND-P001439:locality-cancelled}\textbf{\texttt{GELFAND-0156-U013}} (locality-cancelled).\nopagebreak[4]\par
Before “Main [abbr.], that locality…”, a short word ending in a descending stroke is crossed out.\par
Candidates: Here [source: «Тут»]; Rod [source: «Род»]. Unresolved in the diplomatic source.\par
Historical source review: bbox [285,2750,418,2835], crop [280,2735,760,2850] at 3×. The original strokes overlap the cancellation. The following uncancelled abbreviation is transcribed “Осн.” without expansion.\par
Current source comparison: The 3×/5×/10× views retain the cancelled candidates. The Russian abbreviation remains unexpanded; the English rendering “main” is a translation of the abbreviated sense, not an added source word.\par\medskip
\subsection*{\texttt{GELFAND-P001440}}
\par\noindent\leavevmode\phantomsection\label{GELFAND-0156-U014}\label{GELFAND-P001440:structures-word}\textbf{\texttt{GELFAND-0156-U014}} (structures-word).\nopagebreak[4]\par
At the start of a separate paragraph, before “structures introduced by G. and Dorfman [abbr.]”, an abbreviation is underlined with a wavy line; it resembles “исход.”.\par
Candidates: initial. [uncertain source abbreviation: «исход.»]. Unresolved in the diplomatic source.\par
Historical source review: bbox [324,671,573,748], crop [315,654,581,767] at 5× and line [280,646,2300,871] at 3×. The first and middle signs are joined; the abbreviation was not expanded or replaced by a term from outside knowledge.\par
Current source comparison: The 3×/5×/10× views retain the joined letters and uncertain abbreviation; no external term is substituted.\par\medskip
\par\noindent\leavevmode\phantomsection\label{GELFAND-0156-U015}\label{GELFAND-P001440:laurent-indices}\textbf{\texttt{GELFAND-0156-U015}} (laurent-indices).\nopagebreak[4]\par
In (6), the lower summation index resembles j=−∞, whereas the coefficient subscript and exponent of ξ are written i.\par
Candidates: sum\_\{j=-∞\}\textasciicircum{}N c\_i(z) ξ\textasciicircum{}i; sum\_\{i=-∞\}\textasciicircum{}N c\_i(z) ξ\textasciicircum{}i. Unresolved in the diplomatic source.\par
Historical source review: bbox [937,2065,1345,2258], crop [885,2060,1370,2270] at 4×. The visible j/i mismatch is retained in type. The mathematical expectation of identical indices is recorded here rather than used for a silent correction. C(z)[[ξ]] is also retained as in the source.\par
Current source comparison: The whole series at 3× and the lower index at 5×/10× still favour a j-like descending stroke. No dummy-index normalization is made; the separate consistency note cannot establish the glyph.\par\medskip
\par\noindent\leavevmode\phantomsection\label{GELFAND-0156-U016}\label{GELFAND-P001440:laurent-argument}\textbf{\texttt{GELFAND-0156-U016}} (laurent-argument).\nopagebreak[4]\par
In (6), the sign inside c\_i(…) has been heavily rewritten: z is clearly visible above additional underlying strokes.\par
Candidates: z over ξ; a retraced z. Unresolved in the diplomatic source.\par
Historical source review: bbox [1101,2110,1200,2218], the same crop [885,2060,1370,2270] at 4×. The main text gives the upper z; the earlier sign is not declared deciphered.\par
Current source comparison: At 3×/5×/10×, the upper z remains legible but the overwritten layer is not established; both historical candidates remain.\par\medskip
\subsection*{\texttt{GELFAND-P001441}}
\par\noindent\leavevmode\phantomsection\label{GELFAND-0156-U017}\label{GELFAND-P001441:statement-cancelled}\textbf{\texttt{GELFAND-0156-U017}} (statement-cancelled).\nopagebreak[4]\par
Between the empty bracket with subscript i and “defines”, dense repeated cancellation strokes occur; the original letters cannot reliably be separated.\par
Candidates: none recorded. Unresolved in the diplomatic source.\par
Historical source review: Original 2408×3480; bbox (743,951,1053,1112), magnification 4×. The cancelled content remains illegible and was not reconstructed from the statement of the assertion.\par
Current source comparison: The 3×/5×/10× views do not justify a new reading. The cancelled fragment remains explicitly illegible in both languages.\par\medskip
\par\noindent\leavevmode\phantomsection\label{GELFAND-0156-U018}\label{GELFAND-P001441:final-abbreviation}\textbf{\texttt{GELFAND-0156-U018}} (final-abbreviation).\nopagebreak[4]\par
In the last line, between “quantum” and “currents”, is a short word with a full stop resembling “алг.”.\par
Candidates: alg. [source: «алг.»]; int. [source: «инт.»]. Unresolved in the diplomatic source.\par
Historical source review: Original 2408×3480; bbox (1380,3280,2140,3412), magnification 3×; additionally (1790,3290,2215,3400), 4×. “Токов” (currents) is read from its final в; the preceding abbreviation is not expanded. “Алг.” is preferred from its strokes but remains uncertain.\par
Current source comparison: The abbreviation at 3×/5×/10× still permits uncertainty. “Alg.” is printed provisionally and the alternative “int.” remains in the full note.\par\medskip
\subsection*{\texttt{GELFAND-P001442}}
\par\noindent\leavevmode\phantomsection\label{GELFAND-0156-U019}\label{GELFAND-P001442:kolchin-cancelled}\textbf{\texttt{GELFAND-0156-U019}} (kolchin-cancelled).\nopagebreak[4]\par
Before the clearly written “Kolchin”, a word beginning “Кол…” is crossed out; the visible loops permit a soft sign.\par
Candidates: Kol… [source: «Кол…»]; Kol'… [source: «Коль…»]; Kolchin [source: «Колчин»]. Unresolved in the diplomatic source.\par
Historical source review: Original 2408×3480; bbox (1036,1458,1510,1554), magnification 4×. Complete identity of the cancelled word with the following surname was not established.\par
Current source comparison: The 3× full cancellation and 5×/10× tighter beginning do not prove its complete spelling; the following clear surname is not copied backwards into it.\par\medskip
\par\noindent\leavevmode\phantomsection\label{GELFAND-0156-U020}\label{GELFAND-P001442:heisenberg-sign}\textbf{\texttt{GELFAND-0156-U020}} (heisenberg-sign).\nopagebreak[4]\par
In the right-hand explanation of (21), between “alg.” and “Heisenberg”, at the start of the second line, is a sign with two sloping uprights and transverse strokes; it resembles H/Н or a hash.\par
Candidates: H; N [Cyrillic letter «Н»]; \#; crossed-out Cyrillic letter Н. Unresolved in the diplomatic source.\par
Historical source review: Original 2408×3480; bbox (1812,2047,2268,2239), magnification 3×. It cannot confidently be determined whether the sign names the algebra, begins a cancelled word, or is a separate annotation.\par
Current source comparison: The isolated sign at 3×/5×/10× retains both uprights and crossing strokes but not its function. No algebra designation is inferred from context.\par\medskip
\par\noindent\leavevmode\phantomsection\label{GELFAND-0156-U021}\label{GELFAND-P001442:quadratic-cancelled}\textbf{\texttt{GELFAND-0156-U021}} (quadratic-cancelled).\nopagebreak[4]\par
After “is isomorphic” and before “quadratic algebra [abbr.]”, a cancellation begins with something resembling “ква…”.\par
Candidates: quad… [source: «квад…»]; qua… [source: «ква…»]. Unresolved in the diplomatic source.\par
Historical source review: Original 2408×3480; bbox (1610,2525,1859,2625), magnification 4×. An abandoned first version of the following abbreviation is likely, but the hidden letters were not reconstructed. The isomorphism wording was retained without mathematical correction.\par
Current source comparison: The 3×/5×/10× views do not restore the hidden ending. The written claim about Lie and quadratic algebras remains source-literal; its possible interpretation is discussed separately.\par\medskip
\subsection*{\texttt{GELFAND-P001443}}
\par\noindent\leavevmode\phantomsection\label{GELFAND-0156-U022}\label{GELFAND-P001443:operator-letter}\textbf{\texttt{GELFAND-0156-U022}} (operator-letter).\nopagebreak[4]\par
At the left of the operator definition, before ν/2, is a small sign with an upper bend and a curved lower foot; a similar sign occurs in both exponential factors in (23).\par
Candidates: r; τ. Unresolved in the diplomatic source.\par
Historical source review: Original 2408×3480; bbox (343,843,794,997), magnification 5×; the definition and (23) were also inspected in bbox (194,700,2147,1425), 1.4×. The formulas provisionally use r; R in (25) does not itself prove this letter. The numerator of the fraction is read as ν rather than h.\par
Current source comparison: The tighter 3×/5×/10× sign views do not settle r versus τ. Consistency of the same provisional name across the formulas does not decipher the handwriting.\par\medskip
\par\noindent\leavevmode\phantomsection\label{GELFAND-0156-U023}\label{GELFAND-P001443:boundedness-symbol}\textbf{\texttt{GELFAND-0156-U023}} (boundedness-symbol).\nopagebreak[4]\par
On the right of (23), before (1), is a heavily retraced oval sign with internal strokes and two underlines beneath it.\par
Candidates: O(1); 0(1); O(1) written over an earlier sign. Unresolved in the diplomatic source.\par
Historical source review: Original 2408×3480; bbox (1190,1158,1307,1304), magnification 8×. O(1) agrees with the written limit z→w, but the lower layer of the retracing is not established. The double underline is retained.\par
Current source comparison: The 3×/5×/10× views show the oval and both underlines, but do not establish the earlier layer. No zero or boundedness assertion is proved merely by choosing a candidate reading.\par\medskip
\par\noindent\leavevmode\phantomsection\label{GELFAND-0156-U024}\label{GELFAND-P001443:formula-subline}\textbf{\texttt{GELFAND-0156-U024}} (formula-subline).\nopagebreak[4]\par
Below “Form”, before reference (23), a short word with a full stop has been added; its beginning resembles “ско…” and the middle strokes are faint.\par
Candidates: of the bracket. [source: «скобки.»]; brack. [source: «скоб.»]; of the part. [source: «части.»]. Unresolved in the diplomatic source.\par
Historical source review: Original 2408×3480; bbox (250,2296,694,2464), magnification 4×, and (270,2375,541,2459), 6×. “Скобки.” is preferred, but the faint letters do not permit a final reading; its lower placement is marked in body.tex.\par
Current source comparison: The 3×/5×/10× views retain the below-line position and faint letters. “Of the bracket” remains provisional, with the other full candidates preserved.\par\medskip
\subsection*{\texttt{GELFAND-P001444}}
\par\noindent\leavevmode\phantomsection\label{GELFAND-0156-U025}\label{GELFAND-P001444:reference-name}\textbf{\texttt{GELFAND-0156-U025}} (reference-name).\nopagebreak[4]\par
In the subtitle after “On work [abbr.]”, the surname begins “Пар…” and ends “…сарати”; the group of strokes after the first р permits т or н.\par
Candidates: Partasarati [source: «Партасарати»]; Parnasarati [source: «Парнасарати»]. Unresolved in the diplomatic source.\par
Historical source review: Original 2408×3480; bbox (420,198,825,295), magnification 5×. The surname was not normalized from an external reference; the visible subtitle was not replaced by the catalogue.\par
Current source comparison: The 3×/5×/10× views do not uniquely establish the middle consonant. Both literal surname candidates remain; no external spelling is silently substituted.\par\medskip
\par\noindent\leavevmode\phantomsection\label{GELFAND-0156-U026}\label{GELFAND-P001444:category-cancelled}\textbf{\texttt{GELFAND-0156-U026}} (category-cancelled).\nopagebreak[4]\par
After M(Λ,Ω) and the dash before “category”, several letters are crossed out; a long lower stroke resembling р is visible.\par
Candidates: pr… [source: «пр…»]; cat… [source: «кат…»]. Unresolved in the diplomatic source.\par
Historical source review: Original 2408×3480; bbox (963,653,1110,729), magnification 5×. Dense cancellations prevent restoration of all the letters.\par
Current source comparison: The 3×/5×/10× views do not restore the cancelled word; both beginnings remain candidates only.\par\medskip
\par\noindent\leavevmode\phantomsection\label{GELFAND-0156-U027}\label{GELFAND-P001444:derived-cancelled}\textbf{\texttt{GELFAND-0156-U027}} (derived-cancelled).\nopagebreak[4]\par
Between “such” and “subcategories”, an expression D with a superscript and argument A in parentheses is crossed out. The superscript has a loop resembling e or ℓ.\par
Candidates: D\textasciicircum{}ℓ(A); D\textasciicircum{}e(A); D\textasciicircum{}b(A). Unresolved in the diplomatic source.\par
Historical source review: Original 2408×3480; bbox (933,1344,1247,1492), magnification 5×. This is a cancelled expression, not part of the active formula; the superscript was not corrected to the standard one expected from the meaning.\par
Current source comparison: The 3×/5×/10× views retain the cancellation and looped exponent. No candidate is promoted to active text or normalized to b.\par\medskip
\par\noindent\leavevmode\phantomsection\label{GELFAND-0156-U028}\label{GELFAND-P001444:subcategory-sign}\textbf{\texttt{GELFAND-0156-U028}} (subcategory-sign).\nopagebreak[4]\par
Between A and D(Λ,Ω), an inclusion sign is heavily retraced; inside the rounded outline, two sloping lines form an angle.\par
Candidates: ⊂ over <; ⊂ retraced; ∈ partially rewritten. Unresolved in the diplomatic source.\par
Historical source review: Original 2408×3480; bbox (1751,1355,1884,1472), magnification 7×. The final outline is typeset as ⊂; the lower layer of the correction is not established. The paper ruling was not treated as the middle bar of ∈.\par
Current source comparison: The 3×/5×/10× views preserve the final rounded and underlying angular outlines. The paper ruling does not settle an earlier membership sign.\par\medskip
\par\noindent\leavevmode\phantomsection\label{GELFAND-0156-U029}\label{GELFAND-P001444:derived-exponent}\textbf{\texttt{GELFAND-0156-U029}} (derived-exponent).\nopagebreak[4]\par
In D(A)=D(Λ,Ω), the left D has a superscript with a tall narrow loop and a lower turn to the right; its shape is closer to ℓ.\par
Candidates: ℓ; e; b. Unresolved in the diplomatic source.\par
Historical source review: Original 2408×3480; bbox (483,1471,632,1553), magnification 8×. ℓ is typeset with an explicit caveat. The mathematical expectation of the bounded derived category notation D\textasciicircum{}b was not treated as proof that b is written.\par
Current source comparison: The 3×/5×/10× views retain the narrow loop; ℓ stays provisional. Category-theoretic expectations cannot identify the stroke.\par\medskip
\par\noindent\leavevmode\phantomsection\label{GELFAND-0156-U030}\label{GELFAND-P001444:reflection-quiver}\textbf{\texttt{GELFAND-0156-U030}} (reflection-quiver).\nopagebreak[4]\par
The four-vertex quiver to the left of the reflection definition is labelled Vα, Vβ, Vγ and f. Above the horizontal edge, overlapping strokes resemble g or φ and may be cancelled. The directions of five arrows can be distinguished.\par
Candidates: g above Vα→Vγ; φ above Vα→Vγ; a cancelled label above Vα→Vγ. Unresolved in the diplomatic source.\par
Historical source review: Original 2408×3480; bbox (117,2305,830,2730), magnification 3×. Edges leaving Vα point to the upper, right and lower vertices; the other two point to the right. The lower vertex is unlabelled. Because the label is uncertain, the entire drawing was retained as SourceFigure without an invented recreation.\par
Current source comparison: The historical SourceFigure decision is preserved above. The new partial native draft reproduces the five clear directions, three vertex labels and f beside the exact original crop. The middle label remains explicitly [g/φ ?] after 3×/5×/10× inspection; no label is supplied to the lower vertex. The separate kernel/cokernel calculation cannot identify the sign. PDF layout remains unreviewed.\par\medskip
